\section{Discussion}
\label{sec:discussion}

The studies in Section~\ref{sec:results} show that the usefulness of machine learning, physics-based modelling, and data integration in LIPAS depends on the hazard being forecast and on the purpose of the output. This section interprets those results, states what they do and do not show, and identifies limitations that bound the research-question answers in Section~\ref{sec:conclusion}.

Combining physics and machine learning is not itself new \citep{karpatne2017,willard2022,camporeale2019}. The contribution of LIPAS is narrower: hazard-specific weighting, uncertainty-aware blending, and cross-hazard inputs inside one lunar forecasting and decision-support system. Public heliophysics data enter as inputs to the physics and learned paths rather than as a third modelling method. The results give preliminary support for this design while showing that LIPAS should not replace specialised scientific models or be treated as a fully validated operational forecast.

The April diagnostics used the browser prototype, not the locked hybrid model, and compared it with weakly replicated baselines at a single daytime site. That comparison was designed to expose failure modes rather than to measure skill, and its scores are not pooled with the locked-model results. The temperature overshoot against a Diviner-style dayside relation near peak heating motivated keeping temperature, moonquakes, micrometeorites, and dust on the physics path. Every quantitative claim in Windows~A and~B, the mission simulation, and the ablation therefore rests on the locked production model from the second training run.

The ties on temperature, moonquakes, and dust follow from construction. With locked blend weights $[0.4,0.05,0,0,0,0]$, the hybrid output on those channels is the physics path. The replicated public baselines and evaluation labels for the same channels were built from the same public archives and formulations, so identical RMSE values are expected. They show that unnecessary neural corrections were not forced where physics is already strong; they are not independent validation. Where comparison targets overlap with the construction of LIPAS, agreement partly measures internal consistency.

Solar activity provides the clearest evidence that the learned correction can add information: hybrid RMSE was lower than the replica in both windows (4.10 versus 4.66 in Window~A; 4.99 versus 5.99 in Window~B). The gain is modest and rests on two fourteen-day windows at three sites. Radiation did not improve: the replica had lower RMSE in Window~A (0.102 versus 0.119) and in quiet dayside Window~B (0.00266 versus 0.0188). Radiation should be read as a current weakness of LIPAS, not as a demonstrated strength.

Micrometeorite scores under blend weight zero are circular when the evaluation series carries the same AMS shower term as the physics path. The hybrid RMSE of 0.0 is therefore not an independent win; replica errors measure the absence of that shower term rather than a deficiency of MEM~3 as a whole. The Window~B separation near the Perseids illustrates shower timing on a seasonal background; it does not validate the full NASA micrometeorite modelling framework against measured lunar impacts.

The mission simulation shows why predictive performance and operational usefulness differ. In the radiation-focused EVA scenario, LIPAS issued caution at T+45 and no-go at T+85; the GOES-only path reached caution at T+80 and no-go at T+105; the combined NASA path reached caution at T+55 and the earlier hard no-go at T+70. The earlier LIPAS caution came from upstream solar and geomagnetic drivers, not from detecting a storm that NASA products could not see. The scenario used one injected storm, fixed thresholds, and one hazard. It does not show that LIPAS is safer on every hazard, and a unified system can concentrate risk if its feeds fail.

The ablation study shows cross-hazard sensitivity inside the locked model, not causation. On a relatively quiet portion of Window~A, replacing the solar input raised radiation RMSE from 0.1020 to 0.1050, while the same replacement lowered solar RMSE from 3.899 to 3.290. That solar result is a sensitivity finding for this window; it does not show that live solar data are generally harmful.

Different hazards therefore benefit from different information. Temperature, moonquakes, and dust remain physics-led by design. Solar activity shows a modest learned gain. Radiation does not. Micrometeorite behaviour is driven by explicit shower timing rather than machine learning. LIPAS is not six identical machine-learning models; it assigns influence by hazard.

Several limits bound these conclusions. Continuous co-located surface measurements do not exist for all six hazards, so targets are evaluation labels rather than complete lunar-surface truth. Except for the MEM~3 export, comparisons are author-built replicas of public methods, not runs of agency operational software. The qualitative matrix is not a ranking. The quantitative evidence covers two windows, three sites, and one locked model under relatively quiet radiation conditions.

The strongest conclusion is therefore not that LIPAS replaces specialised NASA or NOAA models. Its value is a common forecasting and decision-support layer in which physics, public heliophysics inputs, cross-hazard context, and uncertainty can be considered together, subject to substantial further validation. Mission planners remain responsible for reviewing and adapting any such outputs before operational use.
